\documentclass[ignorenonframetext, professionalfonts, hyperref={pdftex, unicode}]{beamer}

%Packages to be included

\usepackage{graphicx}
\graphicspath{{../branding/}}

\usepackage{../branding/beamerthemeMLUG}

\usepackage[russian]{babel}
\usepackage[utf8]{inputenc}
\usepackage[T1]{fontenc}

%\usepackage[orientation=landscape, size=custom, width=16, height=9, scale=0.5]{beamerposter}

\usepackage{textcomp}

\usepackage{beamerthemesplit}

\usepackage{ulem}

\usepackage{verbatim}

\usepackage{ucs}


\usepackage{listings}
\lstloadlanguages{bash}

\lstset{escapechar=`,
	extendedchars=false,
	language=sh,
	frame=single,
	tabsize=2, 
	columns=fullflexible, 
%	basicstyle=\scriptsize,
	keywordstyle=\color{blue}, 
	commentstyle=\itshape\color{brown},
%	identifierstyle=\ttfamily, 
	stringstyle=\mdseries\color{green}, 
	showstringspaces=false, 
	numbers=left, 
	numberstyle=\tiny, 
	breaklines=true, 
	inputencoding=utf8,
	keepspaces=true,
	morekeywords={u\_short, u\_char, u\_long, in\_addr}
	}

\definecolor{darkgreen}{cmyk}{0.7, 0, 1, 0.5}

\lstdefinelanguage{diff}
{
    morekeywords={+, -},
    sensitive=false,
    morecomment=[l]{//},
    morecomment=[s]{/*}{*/},
    morecomment=[l][\color{darkgreen}]{+},
    morecomment=[l][\color{red}]{-},
    morestring=[b]",
}

\author[MLUG]{{\bf MLUG}\\Minsk Linux Users Group}

\AtBeginSection[]{%
  \begin{frame}<beamer>
    \frametitle{}
    \tableofcontents[
        sectionstyle=show/shaded, hideallsubsections ]
  \end{frame}
  \addtocounter{framenumber}{-1}% If you don't want them to affect the slide number
}

\AtBeginSubsection[]{%
  \begin{frame}<beamer>
    \frametitle{}
    \tableofcontents[
        sectionstyle=show/hide,
        subsectionstyle=show/shaded/hide, ]
  \end{frame}
  \addtocounter{framenumber}{-1}% If you don't want them to affect the slide number
}


\title{Введение в GNU/Linux}


%%%%%%%%%%%%%%%%%%%%%%%%%%%%%%%%%%%%%%%%%%%%%%%%%
%%%%%%%%%% Begin Document  %%%%%%%%%%%%%%%%%%%%%%
%%%%%%%%%%%%%%%%%%%%%%%%%%%%%%%%%%%%%%%%%%%%%%%%%

\begin{document}

\begin{frame}
	\frametitle{}
	\titlepage
	\vspace{-0.5cm}
	\begin{center}
	%\frontpagelogo
	\end{center}
\end{frame}


\begin{frame}
	\tableofcontents
	[hideallsubsections]
\end{frame}

\section{О курсах}

\mode<all>{
\begin{frame}{Преподаватели}

    \Large{Курс 2018\footnote{В алфавитном порядке}}:

    \begin{columns}
    \column{0.4\textwidth}
        Адамов Юрий
    \column{0.4\textwidth}
        \begin{center}
            \includegraphics[width=0.9\textwidth]{Promwad_logo.png}
        \end{center}
    \end{columns}

    \begin{columns}
    \column{0.4\textwidth}
        Пынькин Денис
    \column{0.4\textwidth}
        \begin{center}
            \includegraphics[width=0.5\textwidth]{COLLABORA_02_RGB.png}
        \end{center}
    \end{columns}
\end{frame}
}

\mode<all>{\begin{frame}{Цель курса}
	\begin{center}
		\Huge
		Увеличение популярности GNU/Linux среди программистов.
	\end{center}
\end{frame}


\begin{frame}{Состав курса}
	\begin{itemize}
		\item Представление об архитектуре GNU/Linux дистрибутива
			\pause
		\item "Ежедневные" навыки работы в консоли
			\pause
		\item Языки программирования
			\begin{itemize}
					\item shell (Bash)
%%					\item Python
%%					\item C
			\end{itemize}
			\pause
		\item Работа с классическими средствами разработки для ОС Linux
			\pause
		\item Все, чего вы не знали и боялись спросить
	\end{itemize}
\end{frame}
}

\section*{GNU/Linux}

\begin{frame}{Основы ОС Linux}

	\begin{block}{Вопрос}
	Почему Linux является самой популярной
	свободной операционной системой?
	\end{block}

	\pause

	\begin{block}{Ответ}
	\begin{itemize}
		\item \textcopyleft -- Copyleft
		\item ``Философия'' Unix
		\item Открытые стандарты
	\end{itemize}
	\end{block}

\end{frame}


%%%%%%%%%%%%%%%%%%%%%%%%%%%%%%%%%%%%%%%%%   
%%%%%%%%%% Content starts here %%%%%%%%%%
%%%%%%%%%%%%%%%%%%%%%%%%%%%%%%%%%%%%%%%%%

\section[Принципы]{Базовые принципы ОС Linux}

\subsection{GNU/Linux}

\mode<all>{\input{../../slides/intro/vocabulary}}

\subsection{Лицензии}

\mode<all>{\begin{frame}{Авторское право и лицензии}

	\begin{block}{Авторское право}
		 Возникает по факту создания ПО 

		\begin{itemize}
			\item Неимущественные права
			\item Имущественные права
		\end{itemize}
	\end{block}

	\pause

    \begin{block}{Лицензии}
		Лицензия -- средство передать какие-либо права на продукт либо его часть.

		Необходима для защиты авторских прав. 
		Средство для возможности законно пресечь несанкционирование копирование,  использование или распространение ПО. 
	\end{block}
\end{frame}


\begin{frame}{Лицензии: открытые и свободные}
	\begin{block}{Р.Столлман: 4 свободы}
		\begin{itemize}
			\item Свобода 0: Свобода запускать программу в любых целях.
			\item Свобода 1: Свобода изучения работы программы и адаптация её к вашим нуждам. 
				Доступ к исходным текстам является необходимым условием.
			\item Свобода 2: Свобода распространять копии,  так что вы можете помочь вашему товарищу.
			\item Свобода 3: Свобода улучшать программу и публиковать ваши улучшения,
				так что всё общество выиграет от этого.
				Доступ к исходным текстам является необходимым условием.
		\end{itemize}
	\end{block}
\end{frame}


\begin{frame}{Лицензии: permissive}
	\begin{columns}
	\column{0.3\textwidth}
		\center\includegraphics[width=2cm]{../../slides/intro/three-arrows@2x}

	\column{0.6\textwidth}

	\begin{itemize}
		\item BSD
		\item MIT
		\item Apache
	\end{itemize}
	\end{columns}

	\begin{block}{I want it simple and permissive.}
		\begin{itemize}
			\item практически не ограничивают свободу действий пользователей ПО и разработчиков, работающих с исходным кодом.
			\item По своему духу, распространение работы под пермиссивной лицензией схоже с помещением работы в общественное
				достояние, но не требует отказа от авторского права.
		\end{itemize}
	\end{block}

\end{frame}


\begin{frame}{\textcopyleft -- Copyleft}

	\begin{columns}
	\column{0.3\textwidth}
		\center\includegraphics[width=2cm]{../../slides/intro/circular@2x}

	\column{0.3\textwidth}

	\begin{block}{Strong}
	\begin{itemize}
		\item GPL v2
        \item GPL v3
		\item AGPL
	\end{itemize}
    \end{block}

	\column{0.3\textwidth}
	\begin{block}{Weak}
	\begin{itemize}
		\item LGPL
		\item MPL
		\item CDDL
	\end{itemize}
    \end{block}

	\end{columns}


	\begin{block}{I care about sharing improvements.}
	
	Авторское лево -- концепция и практика использования законов авторского права для обеспечения 
	невозможности ограничить любому человеку право использовать,  изменять и распространять как 
	исходное произведение,  так и произведения,  производные от него.
	\end{block}


	При копилефте все производные произведения должны распространяться под той же лицензией,
	что и оригинальное произведение.

\end{frame}

\begin{frame}[fragile]{Использование лицензии в проекте}
    \begin{itemize}
        \item Рекомендуется включать файл {\ttfamily LICENSE} в проект
        \item Использовать однозначную отсылку на лицензию в файлах:
            \begin{itemize}
                \item Использовать стандартный заголовок лицензии (если есть)
                \item Использовать SPDX-License-Identifier\\
                    \url{https://spdx.org/licenses/}
                    \begin{lstlisting}[language=C]
/*
 * SPDX-License-Identifier: GPL-2.0
 */
\end{lstlisting}
                \item Использовать URL для определение лицензии
            \end{itemize}
    \end{itemize}
\end{frame}

\begin{frame}[fragile]{Использование \copyright в проекте}
    \begin{block}{Формат записи}
        \begin{verbatim}
<a copyright indicator> <years applicable>, \
        <copyright holder(s)>.
        \end{verbatim}
    \end{block}
    \begin{itemize}
        \item Copyright 1998, Linus Torvalds.
        \item \copyright 2003, 2010, Free Software Foundation, Inc.
        \item Copr. 2005-2012, 2014, John Smith, Jane Doe.
    \end{itemize}
\end{frame}

\begin{frame}[fragile]{Пример}
    \begin{block}{Das U-Boot}
    \begin{lstlisting}[language=C]
/*
* (C) Copyright 2002
* Wolfgang Denk, DENX Software Engineering, wd@denx.de.
*
* (C) Copyright 2010
* Michael Zaidman, Kodak, michael.zaidman@kodak.com
* post_word_{load|store} cleanup.
*
* SPDX-License-Identifier:    GPL-2.0+
*/
\end{lstlisting}
    \end{block}
\end{frame}

\begin{frame}{Полезные ресурсы по лицензиям}
\begin{block}{LFC191}
    В Linux Foundation подготовили бесплатный курс по открытым лицензиям и их применению:\\
    \url{https://training.linuxfoundation.org/training/compliance-basics-for-developers/}
\end{block}
\begin{block}{Подбор лицензии для проекта}
    \url{https://tldrlegal.com/licenses/browse}
\end{block}
\end{frame}

}

\section{Дистрибутивы ОС Linux}

\mode<all>{\input{../../slides/intro/linux-distro}}

\section{Процесс загрузки ОС Linux}

\subsection{Этапы загрузки}

\mode<all>{\input{../../slides/intro/linux-boot}}

\subsection{Ядро Linux}

\mode<all>{\input{../../slides/intro/linux-kernel.tex}}

\subsection{Userspace}

\mode<all>{\input{../../slides/intro/initrd}}

\mode<all>{\input{../../slides/intro/init-process}}

\subsection{Практика}

\mode<all>{\input{../../slides/intro/practice01.tex}}


\mode<all>\begin{frame}[fragile]{Вопросы?}
    \setcounter{tocdepth}{2}
    \tableofcontents

    \bigskip

    \hrulefill
    \begin{columns}
    \column{0.7\textwidth}
            \center{Материалы:}
            \url{https://goo.gl/sEsYPS}
    \column{0.2\textwidth}
        \begin{center}
            \includegraphics[width=0.7\textwidth]{url-qr-2018}
        \end{center}
    \end{columns}

    Исходники: \url{https://github.com/MinskLUG/linux_courses}

\end{frame}

\end{document}
